\naslov{Reševanje polinomskih enačb}

Vemo, da za polinomsko enačbo $x^n + a_{n-1} x^{n-1} + \cdots + a_0 = 0$ obstaja
razširitev polja $\F$, v kateri je enačba razcepna.
Rešitev polinomske enačbe se izraža z radikali, če se da zapisati rešitve enačbe
s pomočjo računskih operacij in korenov.

\begin{definicija}
  Naj bo $K$ polje.
  Razširitvi oblike $K(\sqrt[n]{a})/K$ za $a \in K$ pravimo \pojem{radikalska
	razširitev} polja $K$.
\end{definicija}

Polinomska enačba $p(X) = 0$ je rešljiva z radikali natanko tedaj, ko rešitve
živijo v neki razširitvi $E/F$, pri čemer obstaja končna veriga razširitev $F
\subseteq E_0 \subseteq \ldots \subseteq E_k = E$, kjer so zaporedne razširitve
radikalske.
Problem se prevede na vprašanje iz teorije grup.
Od tu naprej predpostavimo, da so vse razširitve končne.

\begin{definicija}
  Naj bo $E/F$ končna razširitev.
  To je \pojem{normalna razširitev}, če za vsak nerazcepen polinom $p(X) \in
  F[X]$ velja ena od naslednjih možnosti:
  \begin{itemize}
  \item $p$ nima ničle v $E$
  \item $p$ ima vse ničle v $E$
  \end{itemize}
  Ekvivalentno: vsak nerazcepen polinom $p(X) \in F[X]$ z vsaj eno ničlo v $E$
  razpade v linearne faktorje v $E[X]$.
\end{definicija}

\begin{primer}
  $\Q(\sqrt[3]{2})$ nad $\Q$ ni normalna.
  Polinom $p(X) = X^3 - 2$ ne razpade.
\end{primer}

\begin{primer}
  $\Q(\sqrt{2})$ je normalna.
\end{primer}

\vprasanje{Definiraj radikalske in normalne razširitve. Povej primer normalne
  razširitve in razširitve, ki ni normalna.}

\begin{izrek}
  Naj bo $E/F$ končna razširitev.
  Potem je ta razširitev normalna natanko tedaj, ko je $E$ razpadno polje nekega
  polinoma s koeficienti iz $F$.
\end{izrek}

\begin{proof}
  V desno:
  Recimo, da je $E$ normalna.
  Naj bo $p_i(X)$ minimalni polinom $a_i$.
  Definiramo $p(X) = p_1(X) \ldots p_k(X)$.
  Zaradi normalnosti so vse ničle polinoma $p_i(X)$ v $E$, torej so vse ničle
  polinoma $p(X)$ v $E$.
  Velja $F(p) \subseteq E$, ampak $a_1, \ldots, a_k \in F(p)$, torej $F(p)
  \supseteq F(a_1, \ldots, a_k) = E$.

  V levo:
  Recimo, da je $E = F(p)$.
  Vzemimo poljuben nerazcepen $q(X) \in F[X]$, ki ima neko ničlo $a \in E$.
  Naj bo $b$ poljubna ničla $q$.
  Ničli $a$ in $b$ imata isti minimalni polinom, $q$, torej je $F(a) \cong
  F(b)$.
  Izomorfizem $\delta$ lahko razširimo na izomorfizem razpadnih polj $\tau$.
  Predpostavimo lahko, da je $\delta(a) = b$, torej tudi $\tau(a) = b$.
  Ničla $a$ je racionalen izraz, odvisen od $a_1, \ldots, a_k$, torej je tudi
  $b$ racionalen izraz, odvisen od $a_1, \ldots, a_k$, in posledično $b \in E$.
\end{proof}

\vprasanje{Karakteriziraj normalnost za končne razširitve. Dokaži
  karakterizacijo.}

\begin{definicija}
  Polinom $p(X) \in F[X]$ je \pojem{separabilen}, če ima same enostavne ničle.
  Končna razširitev $E/F$ je \pojem{separabilna}, če je minimalni polinom
  vsakega elementa $a \in E$ separabilen.
\end{definicija}

\begin{opomba}
  V karakteristiki $0$ je vsaka končna razširitev separabilna.
\end{opomba}

\begin{izrek}[Primitivni element]
  Vsaka separabilna razširitev je enostavna.
\end{izrek}

\begin{proof}
  Če je $F$ končno polje, je tudi $E$ končno in je $(E^*, \cdot)$ ciklična
  grupa, generirana z $a$.
  Velja $E = F(a_1, \ldots, a_n)$, brez škode za splošnost $a_i \ne 0$, torej
  obstajajo $n_i$, da je $a_i = a^{n_i}$.
  Torej je $E = F(a)$.

  Če pa je $F$ neskončno polje, zapišimo $E = F(a_1, \ldots, a_n)$, ker je
  razširitev končna.
  Poleg tega so $a_i$ algebraični.
  Ker velja $F(a_1, \ldots, a_n) = F(a_1, \ldots, a_{n-2})(a_{n-1}, a_n)$, izrek
  zadošča dokazati za $n = 2$.

  Naj bo torej $E = F(b, c)$ ter $p(X)$ in $q(X)$ minimalna polinoma $b$ in $c$.
  Označimo z $E_1$ razpadno polje polinoma $p(X) q(X)$ in z $b = b_1, \ldots,
  b_r$, $c = c_1, \ldots, c_s$ njune ničle v tem polju.
  Izberimo $\lambda \in F$ tako, da $\lambda \ne (b_j - b){(c - c_k)}^{-1}$ za vse
  $j$ in $k$.
  Tak $\lambda$ res obstaja, ker je polje neskončno.
  Sedaj definiramo $a = b + \lambda c$.
  Očitno je $F(a) \subseteq F(b, c)$; trdimo, da velja tudi vsebovanost v drugi
  smeri.
  Vpeljimo $f(X) = p(a - \lambda X) \in F(a)[X]$.
  Velja $f(c) = p(a - \lambda c) = p(b) = 0$, torej je $c$ hkrati ničla $f(X)$
  in $q(X) \in F[X] \subseteq F(a)[X]$.
  Naj bo $\tilde{q}(X) \in F(a)[X]$ minimalni polinom elementa $c$ nad poljem
  $F(a)$.
  Velja $\tilde{q} | f$ in $\tilde{q} | q$.

  Možne ničle polinoma $\tilde{q}$ so skupne ničle polinomov $q$ in $f$.
  Vemo, da imata skupno ničlo $c$; katere od $c_i$ pa so še ničle $f$?
  Če je $f(c_j) = p(a - \lambda c_j) = 0$, potem je $a - \lambda c_j = b_k$ za
  nek $k$, torej $b + \lambda c - \lambda c_j = b_k$ in posledično $\lambda =
  (b_j - b){(c - c_k)}^{-1}$, kar je po predpostavki nemogoče.
  Torej je $c$ edina ničla $\tilde{q}$.
  Ker je to minimalni polinom in polje $F$ separabilno, je ničla enostavna in
  zato $\tilde{q}(X) = X-c \in F(a)[X]$, torej $c \in F(a)$.
\end{proof}

\begin{primer}
  Razširitev $E = F_2(\sqrt{t})$ je algebraična razširitev, ki ni separabilna.
  Minimalni polinom za $\sqrt{t}$ je $x^2 - t = {(x + \sqrt{t})}^2$.
\end{primer}

\vprasanje{Kdaj je razširitev separabilna? Podaj primer končne razširitve, ki ni
  algebraična. Povej izrek o primitivnem elementu in ga dokaži za končna polja.}

\begin{definicija}
  \pojem{Galoisova grupa razširitve} $E/F$ je grupa
  \[
	\Gal(E/F) = \{\sigma \in \Aut(E) \such \left. \sigma \right|_F =
	\operatorname{id}_F\}.
  \]
\end{definicija}

\begin{lema}
  Naj bo $E/F$ razširitev, $a \in E$ ničla $p(X) \in F[X]$ ter $\sigma \in
  \Gal(E/F)$.
  Potem je $\sigma(a)$ tudi ničla $p(X)$.
\end{lema}

\begin{proof}
  Uporabimo $\sigma$ na zapisu $p(X) = a_n X^n + \cdots + a_0$.
\end{proof}

\begin{definicija}
  Naj bo $p(X) \in F[X]$ polinom.
  \pojem{Galoisova grupa} polinoma $p$ je $\Gal(p) = \Gal(F(p)/F)$.
\end{definicija}

\vprasanje{Definiraj Galoisovo grupo razširitve in polinoma.}

\begin{definicija}
  Normalnim separabilnim razširitvam pravimo \pojem{Galoisove razširitve}.
\end{definicija}

% LocalWords:  separabilno separabilnim
